\section{Conclusion}
\label{sec:conclusion}

Trident is a small robot designed around one task, and the design decisions that matter
most are the ones that removed something. The gripper has no lift axis and no sensor: two
curved jaws on a pair of equal gears cage a coin and slide it, and ridges on the jaws leave
it behind on release. The camera cannot see the floor in front of the robot, and no second
camera was added: the infrared array that follows the line also guides the last part of
every approach, and the floor markings are dimensioned from its pitch so that it can.

Two lessons recur. The first is that a model answers questions that were settled by eye
during the build. The centre of mass lies inside the support triangle, as intended, but the mass
properties of the CAD model put 82\,\% of the weight on the caster, and this, and not the
motors, limits the acceleration of the robot. The simulation showed the
consequence as a skid before the calculation explained it. The second is that the position
of a sensor is a control parameter. The damping of the line follower is set by the distance
of the array ahead of the axle, and an array on the axle would have needed a different
controller.

The simulated mission is complete: the robot follows the line, maps and localises with its
camera, drives around unmapped obstacles, and delivers the three coins by colour, scored by
a referee that it cannot see. The steps that remain are on the hardware: to weigh the robot
on its supports and move the caster and the camera accordingly, to replace the breadboard
wiring of the infrared array, and to run the mission on the printed robot.
